\section{La Consola de Reporte}

Al dar clic en un estudio, entrarás a la pantalla de redacción de reportes, dividida en dos partes.

\subsection{Ajustar el Tamaño de Pantalla (Resizer)}

Puedes ajustar el tamaño del visor o del área de texto según lo necesites:
\begin{enumerate}
    \item La mitad izquierda muestra el \textbf{Visor Nativo PACS} para ver las imágenes. La mitad derecha contiene el formulario de reporte.
    \item Entre ambas partes, hay una barra divisoria vertical delgada en color gris.
    \item Haz clic y mantén presionado sobre esta barra, y \textbf{deslízala hacia la izquierda o derecha} para expandir el panel que necesites. 
    \item Al soltar el clic, el tamaño del visor y del reporte quedarán ajustados a tu comodidad de lectura y escritura.
\end{enumerate}

\begin{figure}[H]
    \centering
    \includegraphics[width=0.8\linewidth]{screenshots/split_pane_editor.png}
    \caption{Ajustando el tamaño del visor PACS.}
\end{figure}

\subsection{Llenado y Envío del Reporte}

En el lado derecho encontrarás los campos necesarios para redactar tu estudio:

\begin{itemize}
    \item \textbf{Descripción del Paciente / Notas de Envío:} Información brindada por el médico tratante o recepción. Puedes modificarla si está incompleta o es necesario aclarar algo.
    \item \textbf{Hallazgos (\textit{Findings}):} Descripción física y anatómica directa de lo que observas en las imágenes (dimensiones, formas, fracturas).
    \item \textbf{Conclusiones Diagnósticas:} Tu diagnóstico médico final sobre las imágenes analizadas.
    \item \textbf{Recomendaciones:} Sugerencias adicionales para el doctor tratante, como solicitar pruebas de ultrasonido extra o pruebas de laboratorio.
\end{itemize}

Para finalizar el reporte:
\begin{enumerate}
    \item Revisa la ortografía de tu texto.
    \item Ve hasta el final de la página y da clic en el botón \textbf{"Generar Reporte" (Generate Report)}.
    \item El sistema agregará tu firma digital al PDF automáticamente y cambiará el estado del estudio a \textsl{"Completado"}.
    \item A partir de este momento, el paciente y los doctores podrán descargar y ver el archivo PDF oficial con tus resultados.
\end{enumerate}
